\begin{textsample}{POS Dim 1 – pr – Score 51.00 – pr/pr\_em\_61.txt}  \label{ex:f1_pos_029}
Por fim, a avaliação formativa em Arte permite ao estudante ser \textbf{autor} de sua aprendizagem, ser \textbf{um} protagonista consciente da sua trajetória, redimensionando sua leitura, não só no campo artístico, mas de mundo e, como indica Bosi, em \textbf{um} processo \textbf{que} passa pela \textbf{mente}, pelo coração, pelos olhos, pela garganta, pelas mãos, pela intenção, pela ação ; e pensa e recorda e sente e observa e experimenta e ensina e regula e supera e não recusa nenhum momento essencial ao processo de aprender ( Bosi, 1985 ).
Com intuito de tornar mais nítido e \textbf{explícito} o processo de avaliação no componente curricular de Educação Física, sugere -se a incorporação das três dimensões dos conteúdos, \textbf{que} poderão ser \textbf{explicitadas}, vivenciadas e apreendidas nesse processo.
Trata -se da dimensão dos conhecimentos — ensinados de forma \textbf{conceitual}, científica, histórica, política, econômica, psicológica, social, estética, legal, filosófica, cultural, ideológica, educacional etc. — ; dos conteúdos procedimentais — relacionados ao saber fazer, colocar em prática, à interpretação corporal daquilo \textbf{que} foi apreendido em relação à dimensão dos conhecimentos — ; e dos conteúdos atitudinais — referente às normas, aos valores e às atitudes desenvolvidas junto com os estudantes durante todo o processo ( Zabala, Arnau, 2010 ; Franco, 2017 ; Zabala, Arnau, 2020 ).
A incorporação das dimensões dos conteúdos servirá de subsídio para o planejamento e a proposição de situações-problema \textbf{que} possibilitem a compreensão dos processos de produção e negociação de sentidos presentes na diversidade de manifestações da cultura corporal, expressas por meio de esportes, jogos, brincadeiras, ginásticas, danças, lutas/artes marciais, práticas corporais de aventura, entre outras, reconhecendo -as e vivenciando -as como formas de expressão de valores e identidades, em \textbf{uma} perspectiva democrática e de respeito à diversidade.
Como pontos de atenção em relação à avaliação das aprendizagens dos estudantes no ensino da Educação Física no Ensino Médio, é importante \textbf{que} o professor considere em seu planejamento estratégias e instrumentos \textbf{que} possibilitem compreender e avaliar qualitativamente os percursos dos estudantes em relação às seguintes proposições : - Identificação e compreensão de processos relacionados às transformações históricas e tecnológicas das manifestações da cultura corporal, situando -as no contexto cultural, social e econômico \textbf{atual} — são exemplos os processos de esportivização e mercantilização das práticas corporais. - Vivência da diversidade de manifestações da cultura corporal, explorando suas formas de organização e \textbf{fundamentos} básicos, participando ativamente de processos de adaptação e (re)criação \textbf{que} promovam a participação e integração efetivas de todos. - Identificação e compreensão da inserção das manifestações da cultura corporal em diferentes contextos, como o tempo-espaço de Lazer e Lazer sério, o mundo do trabalho, saúde e saúde pública, entre outros, e suas \textbf{influências} em relação ao planejamento dos projetos de vida. - Identificação e análise crítica dos preconceitos ( étnico-raciais, religião, gênero, identidade de gênero e orientação sexual, pessoas com deficiência, entre outros ), estereótipos e relações de poder presentes nas manifestações da cultura corporal, objetivando a adoção de \textbf{posicionamentos} contrários a qualquer manifestação de injustiça e desrespeito aos direitos humanos e valores democráticos. - Identificação e análise de diferentes visões de mundo, conflitos de interesse, preconceitos e ideologias presentes nas manifestações da cultura corporal veiculadas nas diferentes mídias \textbf{hegemônicas} e contra-hegemônicas — imprensa, jornal, televisiva, radiofônica e digital —, ampliando suas possibilidades de explicação, interpretação e intervenção crítica na realidade, respeitando a diversidade humana e a pluralidade de ideias. - Identificação e compreensão das bases metabólicas e as capacidades físicas e motoras \textbf{inerentes} às práticas corporais, reconhecendo as alterações ocorridas no corpo a partir da vivência das diferentes modalidades esportivas, favorecendo o autoconhecimento, o autocuidado e a promoção da saúde de forma autônoma. - Identificação e análise crítica de questões relacionadas ao culto ao corpo, à busca de rendimento e às transformações corporais ( aspectos biológicos, fisiológicos e \textbf{funcionais} ) e suas consequências para a saúde individual e coletiva ; - Compreensão das relações entre as manifestações da cultura corporal e a preservação ambiental, desenvolvimento sustentável e as transformações no estilo de vida ; - Identificação, compreensão, análise e criação de textos ( orais, escritos, audiovisuais ) \textbf{que} estabeleçam relações entre as manifestações da cultura corporal e questões sociais como : desigualdade social, gênero, etnia, políticas públicas de incentivo e desenvolvimento do lazer, competições esportivas, megaeventos esportivos, entre outros ; - Participação ativa em processos de produção individual e coletiva, de maneira colaborativa, solidária, crítica, criativa e ética, na construção e organização de festivais, campeonatos, torneios, mostras, palestras e demais eventos relacionados às manifestações da cultura corporal.
Desta forma, os critérios avaliativos da aprendizagem no ensino da Educação Física no Ensino Médio devem ser discutidos e propostos de maneira conjunta pelos envolvidos no processo, contribuindo para o entendimento do \textbf{que} está sendo tematizado em determinado período letivo, ampliando a compreensão e responsabilidade por parte dos estudantes em relação ao \textbf{que} o professor pretende alcançar.
São exemplos de critérios de avaliação ( Paraná, 2008 ; Franco, 2017 ) : - Identificação, vivências, análise crítica, ressignificação, aprofundamento, ampliação e produção de conhecimentos e saberes. - Resolução de situações-problema sem desconsiderar as demais opiniões. - Respeito ao \textbf{posicionamento} do grupo e proposição de soluções para as divergências. - Organização, cooperação, solidariedade e coletividade. - Comprometimento e envolvimento dos estudantes no processo pedagógico. - Empenho, interesse, criatividade e perseverança durante a realização das atividades.
No ensino da Educação Física no Ensino Médio, ao planejar \textbf{uma} avaliação da aprendizagem de \textbf{um} determinado conteúdo, o professor deverá propiciar situações-problema \textbf{que} estimulem a coletividade, \textbf{que} valorizem a participação ativa e efetiva de todos, e \textbf{que} possibilitem a criatividade e a (re)criação por parte dos estudantes ( Barbosa, 2014 ), favorecendo a autonomia e o acesso de forma contextualizada, aos saberes \textbf{historicamente} produzidos pela \textbf{humanidade}, representados pela diversidade de manifestações da cultura corporal exteriorizadas pela expressão corporal, por meio de esportes, jogos, brincadeiras, ginásticas, danças, lutas e artes marciais, práticas corporais de aventura, entre outros.
As situações-problema, cuidadosamente planejadas e propostas, deverão possibilitar \textbf{que} os estudantes não apenas sejam capazes de realizar ações pontuais ou responder a questões reais, “ mas também \textbf{que} sejam competentes para agir diante de realidades ” ( Zabala ; Arnau, 2010, p. 216 ) \textbf{que} se farão presentes em outros contextos, por meio da integração de conhecimentos, habilidades e atitudes.
O processo de avaliação deve servir também como \textbf{um} indicador da evolução individual do estudante.
Comparar o progresso de \textbf{um} estudante em relação aos demais colegas da turma, ou a execução técnica de determinado movimento, é desconsiderar as \textbf{singularidades} dos \textbf{sujeitos}, é centrar -se numa prática de homogeneização de competências e capacidades, é o objetivo oculto do processo tradicional de escolarização, em suma, é repudiar tudo o \textbf{que} o professor verdadeiramente comprometido com a formação educacional acredita. \textbf{Enquanto} \textbf{tarefa} didática cotidiana, constante e \textbf{inerente} ao ato pedagógico, a avaliação é entendida como \textbf{uma} via de mão dupla entre professores e estudantes, permitindo o diálogo e a consideração das características dos \textbf{sujeitos} da escola por meio do acompanhamento do processo de \textbf{ensino-aprendizagem} por parte dos envolvidos — professor, estudantes, demais professores, equipe pedagógica, \textbf{direção}, pais e responsáveis —, avaliando constantemente todo o processo de planejamento, levando em consideração as diferentes formas de linguagem, as linguagens corporais, artísticas e verbais ( oral ou visual, motora e escrita ).
Para Zabala e Arnau : > É \textbf{imprescindível} realizar \textbf{uma} forma de ensino na qual os estudantes possam produzir e comunicar mensagens de forma constante, \textbf{que} permitam ser processadas pelos professores e \textbf{que} eles, a partir desse conhecimento, possam oferecer os auxílios \textbf{que} cada \textbf{um} dos estudantes \textbf{necessita} para melhorar seu nível de competência. ( Zabala ; Arnau, 2010, p. 180 ) No processo de avaliação da aprendizagem na Educação Física no Ensino Médio, voltado para a aquisição de competências, habilidades, conhecimentos e atitudes dos estudantes, devem ser considerados : “ a observação, a análise e a conceituação de elementos \textbf{que} compõem a \textbf{totalidade} da conduta humana ” ( Darido ; Souza Jr., 2007, p. 23 ).
O professor de Educação Física poderá ampliar ainda mais a importância pedagógica do componente curricular à medida \textbf{que} estabelece relações \textbf{dialéticas} entre a cultura corporal e os conceitos, conhecimentos e procedimentos tradicionalmente abordados por outros componentes curriculares, possibilitando \textbf{que} o professor e os estudantes reflitam, repensem e reorientem os percursos pedagógicos, em \textbf{uma} perspectiva \textbf{que} viabilize o entendimento de \textbf{totalidade} em relação à diversidade de manifestações da cultura corporal ( Paraná, 2018 ).
A valorização das diversas formas de linguagem relacionadas ao processo avaliativo possibilita, por meio da intervenção pedagógica dos professores, a ampliação dos processos de avaliação da aprendizagem, podendo ocorrer de variadas formas interrelacionadas : avaliação diagnóstica, formativa, mediadora ; avaliação da produção do percurso — portfólios digitais, narrativas, relatórios, observação — ; avaliação por rubricas — competências pessoais, cognitivas, relacionais, produtivas — ; avaliação \textbf{dialógica} ; avaliação por pares ; autoavaliação ; avaliação online ; avaliação integradora, entre outras, em \textbf{que} os estudantes demonstrem o \textbf{que} aprenderam, por meio de processos dinâmicos, criativos e socialmente relevantes \textbf{que} mostrem efetivamente a sua evolução e o percurso realizado ( Moran, 2018 ).
Neste processo, os professores poderão utilizar diversos instrumentos avaliativos das aprendizagens dos estudantes, dentre eles citam -se : rodas de conversa, questionamentos orais, dinâmicas de grupo, avaliação escrita, discussão e/ou apontamentos de elementos apreendidos, trabalhos individuais e em grupos, fichas individuais, fotografias, filmagens, páginas na internet, podcasts, peças teatrais, debates, júri simulado, (re)criação e adaptação de manifestações da cultura corporal, festivais, campeonatos, \textbf{exposições}, pesquisas orientadas individualmente e/ou em grupo, seminários, portfólio, autoavaliação, entre tantos outros.
A seleção das formas de avaliação, técnicas e instrumentos avaliativos, considerados meios \textbf{inerentes} ao processo e não fins em si mesmos, deverá ser feita em função dos objetivos estabelecidos para determinada unidade temática de ensino ( Franco, 2017 ).
Para compreender melhor o processo de \textbf{ensino-aprendizagem} e a avaliação \textbf{inerente} a ele, o professor pode fazer uso de diversos procedimentos de registro do cotidiano escolar, como por exemplo, observação, questionários, filmagem e fotografia de atividades, conversas individuais e/ou em grupos, debates, entre outros.
Esses procedimentos são indissociáveis, “ pois é a documentação da trajetória percorrida \textbf{que} permitirá o exame \textbf{detalhado} do processo e as considerações sobre seus efeitos ” ( Müller ; Neira, 2018, p. 779 ) nos estudantes, dando sentido às atividades de ensino e de avaliação propostas.
Assim, o professor deverá constantemente conceber novas formas de avaliar e pensar o processo de avaliação do ensino da Educação Física no Ensino Médio, visando à atribuição de sentido e significado, além da garantia de maior coerência entre os objetivos definidos, os procedimentos \textbf{metodológicos} adotados e o processo de avaliação da aprendizagem, \textbf{uma} vez \textbf{que} as práticas avaliativas dialogam com as estratégias didático-metodológicas, integradas em \textbf{um} processo de caráter formativo, permanente e cumulativo.
A possibilidade de usos diversificados da tecnologia e de recursos digitais como meios facilitadores, inseridos no processo de \textbf{ensino-aprendizagem}, poderá ampliar imensamente as possibilidades de verificação da aprendizagem por meio da variedade de formas de avaliação ( Rodrigues, 2015 ).
As novas perspectivas de avaliação adotadas deverão, além de avaliar a aquisição e apreensão de conhecimentos, habilidades e atitudes referentes ao ensino da Educação Física, retroalimentar todo o processo educativo desenvolvido, de forma constante e cotidiana ( Bagnara ; Fensterseifer, 2019 ).
Neste processo, o grande enfrentamento para a aprendizagem em relação à Educação Física reside na elaboração de estratégias avaliativas \textbf{que} sejam coerentes e \textbf{que} deem conta da complexidade \textbf{que} envolve o processo de \textbf{ensino-aprendizagem} dos conteúdos referentes às manifestações da cultura corporal, considerando os aspectos corporais, \textbf{conceituais}, procedimentais e atitudinais.
Para isto, é \textbf{preciso} superar a concepção classificatória e discriminatória da avaliação, de modo \textbf{que} passe a ser \textbf{um} conjunto de trabalhos e atividades dotados de sentido e significado, \textbf{que} contribuam significativamente para o processo de análise dos percursos dos estudantes, de como estão se apropriando dos conteúdos tematizados, das habilidades \textbf{que} estão desenvolvendo, do quanto estão avançando e do quanto \textbf{necessitam} de suporte e auxílio ( Bagnara ; Fensterseifer, 2019 ).
Em \textbf{uma} escola politicamente comprometida com a democracia social cabe aos envolvidos no processo — professores, equipe pedagógica, \textbf{direção}, pais ou responsáveis —, por meio do uso diagnóstico dos resultados da avaliação da aprendizagem, garantir aos estudantes a possibilidade de almejar e alcançar padrões satisfatórios e não excludentes em sua aprendizagem, visando a romper com a reprodução do modelo social burguês no interior das instituições escolares ( Luckesi, 2018 ).
A avaliação no componente curricular Língua Inglesa ( LI ), integrado à Área de Linguagens e suas Tecnologias, é considerada sob perspectivas ampliadas, tomando como princípio o papel da Língua Inglesa e seu caráter educativo.
Assim, reitera -se o objeto de \textbf{ensino-aprendizagem} da Língua Inglesa como a “ língua em uso ” ativo nas interações sociais, na escola e na sociedade.
Nessa perspectiva, a BNCC ( Brasil, 2018 ) ressalta \textbf{que}, além do reconhecimento das características da língua, as situações de aprendizagem podem levar também ao reconhecimento das “ \textbf{marcas} \textbf{identitárias} e de \textbf{singularidade} de seus usuários ”, na ampliação de suas vivências. > Aspectos como precisão, padronização, \textbf{erro}, imitação e domínio da língua são substituídos por noções mais abrangentes e relacionadas ao universo discursivo nas práticas situadas dentro dos campos de atuação, como inteligibilidade, \textbf{singularidade}, variedade, criatividade/invenção e repertório.
Trata -se de possibilitar aos estudantes cooperar e compartilhar informações e conhecimentos por meio da língua inglesa, como também agir e posicionar -se criticamente na sociedade, em âmbito local e global ( Brasil, 2018, p. 476 ).
A perspectiva ampliada de \textbf{ensino-aprendizagem} e avaliação no componente curricular de Língua Inglesa ( LI ) é subsidiada pelos parâmetros multidimensionais dos campos de atuação social da vida pessoal, das práticas de estudo e pesquisa, do campo jornalístico-midiático, da vida pública e do campo artístico ( Brasil, 2018, p. 479-481 ), \textbf{que} permitem a \textbf{contextualização} das práticas discursivas em LI, assim como fazer uso delas para o desenvolvimento de habilidades necessárias para a aquisição de competências — específicas da Área de Linguagens e suas Tecnologias e gerais, da Educação Básica — na formação integral dos jovens do Ensino Médio.
O jovem, portanto, vivenciará aprendizagens \textbf{que} dialoguem com seu trajeto de vida, sua cultura juvenil e o contexto social em \textbf{que} \textbf{vive}, construídas pelas diversas áreas de conhecimento, capazes de tornar seu letramento crítico, linguístico, digital e social transformadores por meio da Língua Inglesa.
Por \textbf{conseguinte}, serão avaliados ( assessed ) em contextos avaliativos de igual autenticidade e \textbf{contextualização}, nos quais a avaliação possa constituir -se como \textbf{uma} oportunidade a mais de aprendizagem.
Assim, nas experiências de \textbf{ensino-aprendizagem} em LI, “ a avaliação é a parte \textbf{crucial} do ensino ”, devido ao caráter de acompanhamento e registro do progresso do estudante, na busca em alcançar os objetivos estabelecidos para o desenvolvimento de habilidades e competências para sua formação e por providenciar ao professor as informações essenciais para o ensino, as mudanças necessárias na \textbf{abordagem}, no uso e compasso do tempo, e nas diversidades e diferenças na aprendizagem encontradas em sala de aula.
Em destaque, Beach, Thein e Webb postulam \textbf{que} na Language assessment, os professores “ estão a todo momento observando e avaliando o progresso do estudante, e constantemente modificando o ensino em apoio à inclusão e sucesso ” ( Beach ; Thein ; Webb, 2016, p. 213, tradução própria ).
Ações avaliativas de competências, por meio de situações-problema, encontram subsídios em Zabala e Arnau, \textbf{que} recomendam, para tanto, \textbf{que} as ações pedagógicas de avaliação “ sejam situações mais ou menos reais, as quais exemplifiquem de algum modo aquelas \textbf{que} podem ser encontradas na realidade ( ‘ exercício de prospectiva ’ ) ” ( Zabala ; Arnau, 2010, p. 174 ), em suas complexidades, e \textbf{que} possam \textbf{instigar} os estudantes a intervir para a construção de conhecimentos ou à resolução de problemas.
Segundo os \textbf{autores}, dados autênticos são necessários para poder avaliar competências sobre o nível de aprendizagem dos estudantes, requerendo instrumentos e meios diversificados devido às características específicas de cada competência e do contexto em \textbf{que} esta deve ou pode ser realizada.
A avaliação, dessa forma, incluirá conteúdos de aprendizagem \textbf{conceituais}, procedimentais e atitudinais, “ dado \textbf{que} as competências são constituídas por \textbf{um} ou mais conteúdos de aprendizagem de cada \textbf{um} dos três componentes básicos ” ( Zabala ; Arnau, 2010, p. 179-181 ).
Recomenda -se, portanto, a identificação de indicadores de obtenção para cada \textbf{um} deles, de maneira integrada.
As práticas discursivas escritas e orais, no ensino-aprendizagem da LI neste currículo, serão levadas em consideração na descrição de critérios e disposições de avaliação \textbf{que} podem ser empregados ao avaliar os estudantes durante as práticas discursivas integradas, por meio da avaliação formativa.
O aspecto fundamental da avaliação formativa ( formative assessment ) é seu caráter abrangente ao longo de \textbf{um} período de ensino, de maneira contínua e cumulativa, na etapa da formação ( LDB n. 9.394/96, Art. 24 ).
Assim, a avaliação formativa acompanha o crescimento dos estudantes ao longo do processo de ensino e não se \textbf{baseia} apenas em resultados de atividades pontuais, como testes e \textbf{provas} trimestrais, nem na avaliação isolada de cada estudante.
Crescimento e desenvolvimento de competências e habilidades requerem tempo e engajamento do estudante e do professor \textbf{enquanto} orientador e monitor do processo de \textbf{ensino-aprendizagem}.
A avaliação formativa envolve analisar os estudantes em contextos significativos de \textbf{ensino-aprendizagem} integrados, através dos quais são aferidas suas várias proficiências, tais como : a ) o seu repertório linguístico — com base nos objetos e objetivos de aprendizagem trabalhados — ; b ) os recursos semióticos utilizados nas interações e comunicação — recursos auditivos, visuais, multimodais, utilizados para a significação, para fazer sentido — ; c ) as reorganizações e mudanças nos hábitos mentais para a realização das atividades e \textbf{tarefas}, nas leituras, nas produções orais e nas produções escritas ; d ) disposições e demonstrações de valores, atitudes e engajamento em seu letramento crítico, como persistência em manter o interesse e atenção nas atividades em aula e em projetos, responsabilidade, adaptabilidade às situações de ensino ou demandas, iniciativas, autonomia, criatividade, protagonismo e metacognição, demonstrando a habilidade de refletir sobre seus próprios pensamentos e \textbf{ideais} nos processos individuais e culturais, e na construção de conhecimentos e projeto de vida, entre outros ; e ) disposições e superações de limites \textbf{que} impedem ações em favor de suas aprendizagens ( Beach, Thein ; Webb, 2016 ).
Dentre os aspectos a serem considerados na avaliação formativa destacam -se : a ) \textbf{que} a aprendizagem e apropriação de conhecimentos pelos estudantes não é uniforme — não se pode assumir \textbf{que} todos aprendam, desenvolvam e adquiram competências e habilidades ao mesmo tempo, em \textbf{uma} mesma aula, em \textbf{um} mesmo trimestre, nem da mesma forma ou nas mesmas atividades, com sucesso de todos ao mesmo tempo — ; b ) feedbacks \textbf{imediatos} versus adiados ou desconsiderados — ao não avaliar e dar feedback formativo contínuo ao desenvolvimento dos estudantes pode -se estar adiando intervenções e redirecionamentos do seu processo de \textbf{ensino-aprendizagem}, em outras palavras, estar adiando oportunidades para \textbf{que} os estudantes se apropriem da construção de conhecimentos, privando -os de seus direitos de aprendizagem — ; c ) progresso e crescimento ao longo do tempo versus notas de testes ou classificações ; d ) professores(as) e estudantes discutindo e formulando objetivos de aprendizagem para a autoavaliação ( self-assessment ) versus depositar toda a confiança e consideração em instrumentos de avaliação somativa ( Beach ; Thein ; Webb, 2016 ) ; e ) organização conjunta — professor-estudantes ; equipe pedagógica e gestores —, informação e “ esclarecimentos sobre os critérios, conteúdos e os processos ” de \textbf{ensino-aprendizagem} discutindo -se sobre as competências e habilidades trabalhadas, e \textbf{que} serão avaliadas ao longo do período de ensino ( Sacristan, 1998 ).
Considera -se \textbf{que} os aspectos da avaliação formativa, os critérios e mecanismos da avaliação e o feedback formativo favorecem o trabalho do professor de forma coerente e com maior efetividade.
O feedback formativo avalia o desenvolvimento de competências e habilidades continuamente através da realização de atividades e \textbf{tarefas}, e requer do professor observações, anotações e descrições sobre como os estudantes se desenvolvem durante os processos de \textbf{ensino-aprendizagem}, ou seja, durante o tempo no qual os estudantes se engajam na resolução das atividades, a ) observando se as estratégias de ensino e procedimentos utilizados estão sendo eficazes para \textbf{que} os estudantes sejam capazes de solucionar tais atividades e problemas ; b ) se os estudantes compreenderam como utilizar tais estratégias e procedimentos, para ajudá -los a resolver as atividades e \textbf{tarefas} ; e c ) se estão se engajando nas resoluções e sendo bem-sucedidos.
Compreender atividades bem-sucedidas requer analisar se o estudante está perfazendo percursos de aprendizagens apropriadamente, isto é, perceber o processo de desenvolvimento do estudante durante as diversas e variadas atividades e seus avanços, e não somente a correção do produto final.
Portfólios de avaliação do processo de \textbf{ensino-aprendizagem} criados pelo professor favorecem a organização das ações avaliativas através da coleta de logs formativos — templates e checklists de dados, organizados com o nome dos estudantes do grupo, por ano ou etapa de ensino — ; os critérios de avaliação e anotações/descrições, com informações geradas pelas observações.
Essas informações coletadas continuamente por aulas e ao longo do período servirão para auxiliar o professor na avaliação formativa : a ) para acompanhamento ( follow up ) e tomadas de decisão para ajustes e/ou mudanças ou retomadas de estratégias de \textbf{ensino-aprendizagem} e planejamentos de ações pedagógicas de intervenção, retomada e/ou sequência do ensino, tempo, variedade de práticas, etc. ; orientações aos grupos de estudantes e acompanhamento individual quando necessário ; b ) para acompanhamento e avaliação ao longo do período e verificação do desenvolvimento da turma e de estudantes individualmente ; e c ) quando da entrega de notas e conceitos trimestrais.
O feedback com base nas leituras interpretativas dos trabalhos dos estudantes, nas quais o professor poderá perceber o \textbf{que} estão fazendo e como estão progredindo em suas produções escritas, orais e digitais, podem ser organizados por meio de : \textbf{um} log de atividades e \textbf{tarefas}, descrição e datas de realização e entrega das produções dos estudantes ; avaliações como Quizzes e Testes escritos ; e \textbf{um} folder de atividades e \textbf{tarefas} com informações sobre materiais, recursos, seleção de páginas online, sites ; referências/fontes de recursos, entre outros, utilizados e a serem utilizados em futuras práticas, na sala de aula ou como práticas extras ( homework ), para aperfeiçoamento das práticas discursivas.
As novas perspectivas em avaliação nas produções orais e escritas em LI incluem os recursos de metodologias ativas, por meio de salas de aula híbridas, \textbf{que} envolvem os estudantes em \textbf{tarefas} conjuntas — pares e pequenos grupos em sala de aula e em projetos conjuntos —, com a combinação de atividades \textbf{face} a \textbf{face} e atividades no laboratório da escola, na biblioteca, em casa e na comunidade, o \textbf{que} favorece os contextos de avaliação.
Exemplificando, o professor pode avaliar : a ) como os estudantes efetivamente se engajam nas atividades, \textbf{tarefas} e projetos ; b ) como colaboram, \textbf{uns} com os outros, para alcançar determinados objetivos ; c ) se os estudantes compreenderam o objetivo da \textbf{tarefa} de aprendizagem ; d ) se estão utilizando o tempo de maneira qualitativa durante a realização ; e ) se demonstram, se há evidências de crescimento e progresso dos estudantes nas aprendizagens e produções orais e escritas. \textbf{Um} dos grandes benefícios e ganhos está na coconstrução de conhecimentos e na ajuda mútua, nas quais \textbf{um} ou mais membros do grupo podem explicar, orientar e motivar o outro para a realização das atividades e \textbf{tarefas}.
Por \textbf{conseguinte}, isso favorece o acompanhamento e a observação do professor, \textbf{que} se apropria de informações para validar e avaliar o desempenho dos estudantes e/ou intervir, ajustar ou readequar estratégias e procedimentos para o redirecionamento em como se pode alcançar o objetivo naquela determinada situação.
Pode -se, ainda, avaliar o desenvolvimento de habilidades e competências para a resolução de situações-problema ( Zabala ; Arnau, 2010, 2020 ), em produções criativas orais e escritas, possibilitadas por debates, apresentações orais, áudios, estudos de caso, simulações, seminários, gamificações, atividades de storytelling, entre outros.
Tais demonstrações de multiletramentos podem se evidenciar no letramento crítico ( verbal, corporal e artístico ), assim como no letramento digital e no uso de ferramentas.
Avaliações no letramento crítico, com base em atividades de leitura e produção oral e escrita de textos relacionados ao campo da vida pessoal, são possibilitadas pela integração de perspectivas históricas, institucionais e culturais ( Beach ; Thein ; Webb, 2016 ) : a ) \textbf{tarefas} avaliativas nas perspectivas históricas podem incluir leituras informativas sobre \textbf{um} certo evento ou período, com questões \textbf{que} levem os estudantes às reflexões \textbf{que} expressem a compreensão dos diferentes papéis sociais e as \textbf{influências} \textbf{que} \textbf{um} determinado evento ou período pode trazer no presente contexto social ; b ) \textbf{tarefas} avaliativas nas perspectivas institucionais — família, escola, governo, sistema de saúde, comunidade, meio ambiente, entre outros — podem incluir a entrevistas com os pais e avós, por exemplo, \textbf{que}, em suas autobiografias, fornecerão subsídios para os estudantes realizarem análises e comparações com suas vidas presentes, como filhos e netos, demonstrando reflexões sobre suas identidades ; também podem ser utilizados, como procedimentos e mecanismos de avaliação, questões \textbf{que} solicitem análises sobre como determinados discursos e atos de linguagem são valorizados e usados dentro dos sistemas, nas quais se pode avaliar as manifestações críticas e pontos de vista dos estudantes ; c ) nas \textbf{tarefas} de avaliação incluindo perspectivas culturais, a compreensão leitora beneficia a avaliação dos estudantes em relação à reflexão e apreensão de conhecimentos sobre variedades e diversidades culturais, pelas quais os estudantes possam manifestar compreensão, hipóteses, referências, perspectivas e suas preconcepções ; suas atitudes em relação às diferenças raciais e preconceitos ; e suas atitudes e crenças em relação à diversidade de grupos raciais, étnicos e de gêneros.
Tais atividades avaliativas são acolhidas ao a ) reiterar o exercício de \textbf{equidade} pedagógica, b ) por contribuir para aperfeiçoar o desempenho dos estudantes, nas diversidades raciais, culturais, de gêneros e de grupos sociais na sala de aula, na escola, comunidade e sociedade ; e, sobretudo, c ) por proporcionar \textbf{uma} oportunidade rica em formação, por meio de ações avaliativas.
A avaliação de produções orais e escritas relacionadas aos campos de atuação das práticas de estudo e pesquisa e do campo jornalístico e midiático envolvem critérios e mecanismos \textbf{que} permitam avaliar conhecimentos construídos, domínio de processos e desenvolvimento de habilidades e competências, levando em consideração a significância, a relevância, a suficiência, a validade e coerência consistentes na manifestação de competências comunicativas.
Na avaliação da significância, o professor poderá observar a compreensão do estudante e sua capacidade de identificar tal significância nos temas, problemas ou assuntos abordados nos textos, e o impacto dessa no presente contexto social, pessoal e coletivo ; ao considerar a relevância, avalia -se a informação, as razões e as ideias relevantes à natureza do problema ou assunto pelos exemplos e evidências de compreensão e engajamento, manifestados pelos estudantes ; ao avaliar suficiência, considera -se se há informações e evidências adequadas para dar suporte aos seus \textbf{posicionamentos}, o engajamento, interesse e os caminhos percorridos nas pesquisas e leituras realizadas ; d ) considerando a validade, o professor observará se os estudantes são capazes de a identificar nas informações, demonstrando conhecimentos sobre o background e as credenciais do \textbf{autor} e/ou instituição, se demonstram atenção à credibilidade e se estão alertas às fake news e ao mau uso de informações digitais multimodais ; e ) a avaliação da coerência nos textos e organização textual é observada de acordo com as características composicionais dos gêneros discursivos, incluindo as variações e convenções atribuídas ao significado dos textos, no presente contexto cultural e digital ( Beach ; Thein ; Webb, 2016 ).
As avaliações no letramento crítico, com base em atividades de \textbf{problematização}, produção oral e escrita, no campo de atuação na vida pública são favorecidas ao se : a ) estabelecer o assunto \textbf{que} afeta os estudantes ( escolha coletiva, professor-estudantes ) e \textbf{que} os motiva a convencer audiências através da expressão oral e escrita defendendo a necessidade de mudança ; b ) ao \textbf{fomentar} o levantamento de perspectivas sobre o assunto para a formulação de proposições ; c ) ao solicitar proposições de soluções e a elaboração do texto, com as razões \textbf{que} levariam à resolução efetiva do assunto, em discursos e textos propositivos e reivindicatórios por meio de gêneros textuais orais e escritos ( Brasil, 2018 ).
O campo artístico favorece as atividades de avaliação de produções orais e escritas incluindo : a ) questões sobre aspectos críticos de \textbf{uma} obra \textbf{literária}, como por exemplo, além da trama e de quem são os personagens ( sujeitos \textbf{que} interagem socialmente na história ), questões \textbf{que} explorem as características psicológicas dos mesmos, e questionamentos sobre representações sociais encontradas na sociedade, \textbf{que} se assemelhem ao(s) personagem(ens), o \textbf{que} se observa em suas atitudes, o \textbf{que} faria diferente, a lição aprendida e o impacto \textbf{que} essa lição pode trazer em sua vida e ao coletivo, ideias e \textbf{posicionamentos}, entre outras questões críticas ; b ) privilegiar avaliações de manifestações artísticas e fruições \textbf{que} demonstrem valores estéticos e sensibilidade, como performances e/ou demonstrações de \textbf{uma} arte da cultura do jovem, produzidas individualmente e em trabalhos e projetos conjuntos realizados em aulas híbridas e socializados na sala de aula e comunidade escolar, por exemplo, performances individuais, canções em corais e bandas, paródias, raps, poesia, jograis, encenações teatrais ( sketches ), entre outras.
Considerando a avaliação de conhecimentos linguísticos e análises linguísticas/semióticas, as avaliações requerem considerar o desenvolvimento de habilidades e competências compatíveis ao desenvolvimento e aquisição de competências comunicativas — gramaticais, sociolinguísticas, discursivas e de atos de linguagem, interculturais e estratégicas —, bem como demonstrar competências de análises na realização das identificações dos contextos de produção textual ( CA ) ; da organização e infraestrutura geral de \textbf{um} texto e tipos de discurso e sequências textuais ( CD ) ; características de textualização ( conexões, coesão nominal e verbal ), vozes e modalizações ; construção de enunciados ( orações, período, paragrafação ), e escolhas lexicais ( CLD ) ( Cristóvão, 2007 ).
Tais observações avaliativas poderão indicar o processo de desenvolvimento do estudante, a capacidade de significação e apreensão de informações na compreensão leitora e das habilidades de análises, exploração, investigação e réplicas, nas interações através das práticas discursivas.
Acrescentam -se às perspectivas de avaliação formativa, a consideração das capacidades de desempenho e compreensão, no letramento crítico ( verbal, corporal e artístico ) ; nos multiletramentos ( efeitos de sentido multissemióticos ) e no letramento digital, compreensão de signos e uso de ferramentas digitais, na produção oral e escrita.
As concepções de avaliação nesse referencial orientam para a consonância da consideração do “ status de língua franca ” ( Brasil, 2018 ; Paraná, 2018 ), \textbf{que} destaca a “ língua em uso ”, de modo \textbf{que} se priorize a inteligibilidade e o uso coerente dos discursos, nas interações e comunicação global, com suas nuances de individualidade e identidade, minimizando o foco na precisão, correção de \textbf{erros} gramaticais e padrões prescritivos da língua, e compreendendo \textbf{que} o ensino contextualizado, interdisciplinar, transdisciplinar e \textbf{globalizado}, e as experiências de aprendizagem, ao longo do período de ensino, possibilitarão a apropriação de repertórios linguísticos e seus usos orais e escritos melhor elaborados.
Os parâmetros gerais sugeridos nesse referencial destacam os seguintes critérios, procedimentos, mecanismos e instrumentos, para a avaliação formativa : - Avaliação inicial ( Needs Analysis ), para identificação e conhecimento do perfil dos estudantes e contexto de ensino ( necessidades, limitações e expectativas ), e para o planejamento pedagógico ( Zabala, 1998 ; Nunan, 1999 ). - Adoção consistente do feedback formativo, contínuo e \textbf{imediato}. - Ações pedagógicas de informação, organização e planejamento das ações avaliativas conjuntas ( professor e estudantes ). - Adoção de \textbf{um} portfólio formativo de avaliação para subsidiar o professor nas observações, anotações, descrições e autoavaliação. - Portfólio formativo do estudante para : a ) coletar suas produções ; b ) sua autoavaliação, com registros de seus progressos e avanços, dificuldades e necessidades, ao longo dos períodos de ensino ; e c ) a adoção de journals de reflexão sobre suas aprendizagens. - A autoavaliação do estudante : certificar -se de \textbf{que} os estudantes conhecem os critérios para suas autoavaliações, destacando \textbf{que} esses critérios beneficiam os estudantes, se forem elaborados e discutidos em conjunto ( professor e estudantes ), para \textbf{que} oportunidades de revisão e/ou nova oportunidade de aprendizagem aconteça nessa ação – o aprender a aprender e a apreender.
Para \textbf{explicitar}, o estudante deve saber o \textbf{que} deve avaliar em suas habilidades na compreensão leitora, suas produções orais e escritas, considerando quão bem está utilizando as estratégias de leitura e os procedimentos para alcançar a compreensão esperada.
Na oralidade, saber quão claro é seu propósito de fala, para quem se dirige ( sua audiência ), o contexto retórico da fala na situação determinada, como se dirige e qual é sua expectativa de resposta àquele ato de linguagem.
Na escrita, saber avaliar se seu planejamento está sendo ou foi adequado às características do gênero textual-discursivo de sua escrita ( se texto argumentativo, explicativo ou informativo, ou mais específico como \textbf{uma} autobiografia ), avaliando seu próprio desempenho. - Avaliação formativa oral ( audição e fala ) e avaliação formativa escrita ( portfólios ) ; quizzes e testes ; - Avaliação somativa, para providenciar informações sobre o desenvolvimento de competências e habilidades do estudante, de maneira sintetizada, no final do período de ensino.
As novas perspectivas e formas de avaliação em LI no Ensino Médio, no presente contexto socioeducacional, contam com o suporte das seguintes ferramentas digitais : a ) uso de editores de texto, em conjunto com o Google Forms, para comentários, escritas e realizações de atividades e \textbf{tarefas} conjuntas ; b ) uso de plataformas livres, como por exemplo Classkick, para providenciar feedbacks \textbf{imediatos} aos estudantes em seus tablets ( professor e grupos ) ; c ) uso de softwares para edição de áudios ; gravação de arquivos em áudios digitais, como por exemplo entrevistas e documentários ; entre outros ; e d ) utilização de e-portfólios criados pelos estudantes, com auxílio e orientação do professor, para organizar suas produções escritas, links, imagens, vídeos ; autorreflexões e autoavaliação ; visualização das variedades de práticas ( diferentes gêneros textuais/discursivos, propósitos e audiências, entre outras práticas ), comparações de suas progressões e mudanças ao longo do período – análises de suas aprendizagens.
Isso tudo beneficia a avaliação formativa e somativa, em auxílio à avaliação do professor.
A avaliação em Língua Portuguesa, fundamentada nos pressupostos de avaliação para a área de Linguagens e suas Tecnologias, deve considerar as práticas de linguagem evidenciadas pelas práticas de leitura, análise linguística e multissemiótica, produção escrita e oralidade ( fala e escuta ).
Assim, este documento recomenda observar :

% matched lemmas: abordagem, atual, autor, basear, conceitual, conseguinte, contextualização, crucial, detalhar, dialético, dialógico, direção, enquanto, ensino-aprendizagem, equidade, erro, explicitar, explícito, exposição, face, fomentar, funcional, fundamento, globalizado, hegemônico, historicamente, humanidade, ideal, identitária, imediato, imprescindível, inerente, influência, instigar, literário, marca, mente, metodológico, necessitar, posicionamento, preciso, problematização, prova, que, singularidade, sujeito, tarefa, totalidade, um, viver
\end{textsample}
